%% bare_jrnl.tex
%% V1.4b
%% 2015/08/26
%% by Michael Shell
%% see http://www.michaelshell.org/
%% for current contact information.
%%
%% This is a skeleton file demonstrating the use of IEEEtran.cls
%% (requires IEEEtran.cls version 1.8b or later) with an IEEE
%% journal paper.
%%
%% Support sites:
%% http://www.michaelshell.org/tex/ieeetran/
%% http://www.ctan.org/pkg/ieeetran
%% and
%% http://www.ieee.org/

%%*************************************************************************
%% Legal Notice:
%% This code is offered as-is without any warranty either expressed or
%% implied; without even the implied warranty of MERCHANTABILITY or
%% FITNESS FOR A PARTICULAR PURPOSE! 
%% User assumes all risk.
%% In no event shall the IEEE or any contributor to this code be liable for
%% any damages or losses, including, but not limited to, incidental,
%% consequential, or any other damages, resulting from the use or misuse
%% of any information contained here.
%%
%% All comments are the opinions of their respective authors and are not
%% necessarily endorsed by the IEEE.
%%
%% This work is distributed under the LaTeX Project Public License (LPPL)
%% ( http://www.latex-project.org/ ) version 1.3, and may be freely used,
%% distributed and modified. A copy of the LPPL, version 1.3, is included
%% in the base LaTeX documentation of all distributions of LaTeX released
%% 2003/12/01 or later.
%% Retain all contribution notices and credits.
%% ** Modified files should be clearly indicated as such, including  **
%% ** renaming them and changing author support contact information. **
%%*************************************************************************


% *** Authors should verify (and, if needed, correct) their LaTeX system  ***
% *** with the testflow diagnostic prior to trusting their LaTeX platform ***
% *** with production work. The IEEE's font choices and paper sizes can   ***
% *** trigger bugs that do not appear when using other class files.       ***                          ***
% The testflow support page is at:
% http://www.michaelshell.org/tex/testflow/



\documentclass[journal]{IEEEtran}
\usepackage{amsmath,amssymb}
\usepackage{booktabs}
\usepackage{graphicx}
%
% If IEEEtran.cls has not been installed into the LaTeX system files,
% manually specify the path to it like:
% \documentclass[journal]{../sty/IEEEtran}





% Some very useful LaTeX packages include:
% (uncomment the ones you want to load)


% *** MISC UTILITY PACKAGES ***
%
%\usepackage{ifpdf}
% Heiko Oberdiek's ifpdf.sty is very useful if you need conditional
% compilation based on whether the output is pdf or dvi.
% usage:
% \ifpdf
%   % pdf code
% \else
%   % dvi code
% \fi
% The latest version of ifpdf.sty can be obtained from:
% http://www.ctan.org/pkg/ifpdf
% Also, note that IEEEtran.cls V1.7 and later provides a builtin
% \ifCLASSINFOpdf conditional that works the same way.
% When switching from latex to pdflatex and vice-versa, the compiler may
% have to be run twice to clear warning/error messages.






% *** CITATION PACKAGES ***
%
%\usepackage{cite}
% cite.sty was written by Donald Arseneau
% V1.6 and later of IEEEtran pre-defines the format of the cite.sty package
% \cite{} output to follow that of the IEEE. Loading the cite package will
% result in citation numbers being automatically sorted and properly
% "compressed/ranged". e.g., [1], [9], [2], [7], [5], [6] without using
% cite.sty will become [1], [2], [5]--[7], [9] using cite.sty. cite.sty's
% \cite will automatically add leading space, if needed. Use cite.sty's
% noadjust option (cite.sty V3.8 and later) if you want to turn this off
% such as if a citation ever needs to be enclosed in parenthesis.
% cite.sty is already installed on most LaTeX systems. Be sure and use
% version 5.0 (2009-03-20) and later if using hyperref.sty.
% The latest version can be obtained at:
% http://www.ctan.org/pkg/cite
% The documentation is contained in the cite.sty file itself.






% *** GRAPHICS RELATED PACKAGES ***
%
\ifCLASSINFOpdf
  % \usepackage[pdftex]{graphicx}
  % declare the path(s) where your graphic files are
  % \graphicspath{{../pdf/}{../jpeg/}}
  % and their extensions so you won't have to specify these with
  % every instance of \includegraphics
  % \DeclareGraphicsExtensions{.pdf,.jpeg,.png}
\else
  % or other class option (dvipsone, dvipdf, if not using dvips). graphicx
  % will default to the driver specified in the system graphics.cfg if no
  % driver is specified.
  % \usepackage[dvips]{graphicx}
  % declare the path(s) where your graphic files are
  % \graphicspath{{../eps/}}
  % and their extensions so you won't have to specify these with
  % every instance of \includegraphics
  % \DeclareGraphicsExtensions{.eps}
\fi
% graphicx was written by David Carlisle and Sebastian Rahtz. It is
% required if you want graphics, photos, etc. graphicx.sty is already
% installed on most LaTeX systems. The latest version and documentation
% can be obtained at: 
% http://www.ctan.org/pkg/graphicx
% Another good source of documentation is "Using Imported Graphics in
% LaTeX2e" by Keith Reckdahl which can be found at:
% http://www.ctan.org/pkg/epslatex
%
% latex, and pdflatex in dvi mode, support graphics in encapsulated
% postscript (.eps) format. pdflatex in pdf mode supports graphics
% in .pdf, .jpeg, .png and .mps (metapost) formats. Users should ensure
% that all non-photo figures use a vector format (.eps, .pdf, .mps) and
% not a bitmapped formats (.jpeg, .png). The IEEE frowns on bitmapped formats
% which can result in "jaggedy"/blurry rendering of lines and letters as
% well as large increases in file sizes.
%
% You can find documentation about the pdfTeX application at:
% http://www.tug.org/applications/pdftex





% *** MATH PACKAGES ***
%
%\usepackage{amsmath}
% A popular package from the American Mathematical Society that provides
% many useful and powerful commands for dealing with mathematics.
%
% Note that the amsmath package sets \interdisplaylinepenalty to 10000
% thus preventing page breaks from occurring within multiline equations. Use:
%\interdisplaylinepenalty=2500
% after loading amsmath to restore such page breaks as IEEEtran.cls normally
% does. amsmath.sty is already installed on most LaTeX systems. The latest
% version and documentation can be obtained at:
% http://www.ctan.org/pkg/amsmath





% *** SPECIALIZED LIST PACKAGES ***
%
%\usepackage{algorithmic}
% algorithmic.sty was written by Peter Williams and Rogerio Brito.
% This package provides an algorithmic environment fo describing algorithms.
% You can use the algorithmic environment in-text or within a figure
% environment to provide for a floating algorithm. Do NOT use the algorithm
% floating environment provided by algorithm.sty (by the same authors) or
% algorithm2e.sty (by Christophe Fiorio) as the IEEE does not use dedicated
% algorithm float types and packages that provide these will not provide
% correct IEEE style captions. The latest version and documentation of
% algorithmic.sty can be obtained at:
% http://www.ctan.org/pkg/algorithms
% Also of interest may be the (relatively newer and more customizable)
% algorithmicx.sty package by Szasz Janos:
% http://www.ctan.org/pkg/algorithmicx




% *** ALIGNMENT PACKAGES ***
%
%\usepackage{array}
% Frank Mittelbach's and David Carlisle's array.sty patches and improves
% the standard LaTeX2e array and tabular environments to provide better
% appearance and additional user controls. As the default LaTeX2e table
% generation code is lacking to the point of almost being broken with
% respect to the quality of the end results, all users are strongly
% advised to use an enhanced (at the very least that provided by array.sty)
% set of table tools. array.sty is already installed on most systems. The
% latest version and documentation can be obtained at:
% http://www.ctan.org/pkg/array


% IEEEtran contains the IEEEeqnarray family of commands that can be used to
% generate multiline equations as well as matrices, tables, etc., of high
% quality.




% *** SUBFIGURE PACKAGES ***
%\ifCLASSOPTIONcompsoc
%  \usepackage[caption=false,font=normalsize,labelfont=sf,textfont=sf]{subfig}
%\else
%  \usepackage[caption=false,font=footnotesize]{subfig}
%\fi
% subfig.sty, written by Steven Douglas Cochran, is the modern replacement
% for subfigure.sty, the latter of which is no longer maintained and is
% incompatible with some LaTeX packages including fixltx2e. However,
% subfig.sty requires and automatically loads Axel Sommerfeldt's caption.sty
% which will override IEEEtran.cls' handling of captions and this will result
% in non-IEEE style figure/table captions. To prevent this problem, be sure
% and invoke subfig.sty's "caption=false" package option (available since
% subfig.sty version 1.3, 2005/06/28) as this is will preserve IEEEtran.cls
% handling of captions.
% Note that the Computer Society format requires a larger sans serif font
% than the serif footnote size font used in traditional IEEE formatting
% and thus the need to invoke different subfig.sty package options depending
% on whether compsoc mode has been enabled.
%
% The latest version and documentation of subfig.sty can be obtained at:
% http://www.ctan.org/pkg/subfig




% *** FLOAT PACKAGES ***
%
%\usepackage{fixltx2e}
% fixltx2e, the successor to the earlier fix2col.sty, was written by
% Frank Mittelbach and David Carlisle. This package corrects a few problems
% in the LaTeX2e kernel, the most notable of which is that in current
% LaTeX2e releases, the ordering of single and double column floats is not
% guaranteed to be preserved. Thus, an unpatched LaTeX2e can allow a
% single column figure to be placed prior to an earlier double column
% figure.
% Be aware that LaTeX2e kernels dated 2015 and later have fixltx2e.sty's
% corrections already built into the system in which case a warning will
% be issued if an attempt is made to load fixltx2e.sty as it is no longer
% needed.
% The latest version and documentation can be found at:
% http://www.ctan.org/pkg/fixltx2e


%\usepackage{stfloats}
% stfloats.sty was written by Sigitas Tolusis. This package gives LaTeX2e
% the ability to do double column floats at the bottom of the page as well
% as the top. (e.g., "\begin{figure*}[!b]" is not normally possible in
% LaTeX2e). It also provides a command:
%\fnbelowfloat
% to enable the placement of footnotes below bottom floats (the standard
% LaTeX2e kernel puts them above bottom floats). This is an invasive package
% which rewrites many portions of the LaTeX2e float routines. It may not work
% with other packages that modify the LaTeX2e float routines. The latest
% version and documentation can be obtained at:
% http://www.ctan.org/pkg/stfloats
% Do not use the stfloats baselinefloat ability as the IEEE does not allow
% \baselineskip to stretch. Authors submitting work to the IEEE should note
% that the IEEE rarely uses double column equations and that authors should try
% to avoid such use. Do not be tempted to use the cuted.sty or midfloat.sty
% packages (also by Sigitas Tolusis) as the IEEE does not format its papers in
% such ways.
% Do not attempt to use stfloats with fixltx2e as they are incompatible.
% Instead, use Morten Hogholm'a dblfloatfix which combines the features
% of both fixltx2e and stfloats:
%
% \usepackage{dblfloatfix}
% The latest version can be found at:
% http://www.ctan.org/pkg/dblfloatfix




%\ifCLASSOPTIONcaptionsoff
%  \usepackage[nomarkers]{endfloat}
% \let\MYoriglatexcaption\caption
% \renewcommand{\caption}[2][\relax]{\MYoriglatexcaption[#2]{#2}}
%\fi
% endfloat.sty was written by James Darrell McCauley, Jeff Goldberg and 
% Axel Sommerfeldt. This package may be useful when used in conjunction with 
% IEEEtran.cls'  captionsoff option. Some IEEE journals/societies require that
% submissions have lists of figures/tables at the end of the paper and that
% figures/tables without any captions are placed on a page by themselves at
% the end of the document. If needed, the draftcls IEEEtran class option or
% \CLASSINPUTbaselinestretch interface can be used to increase the line
% spacing as well. Be sure and use the nomarkers option of endfloat to
% prevent endfloat from "marking" where the figures would have been placed
% in the text. The two hack lines of code above are a slight modification of
% that suggested by in the endfloat docs (section 8.4.1) to ensure that
% the full captions always appear in the list of figures/tables - even if
% the user used the short optional argument of \caption[]{}.
% IEEE papers do not typically make use of \caption[]'s optional argument,
% so this should not be an issue. A similar trick can be used to disable
% captions of packages such as subfig.sty that lack options to turn off
% the subcaptions:
% For subfig.sty:
% \let\MYorigsubfloat\subfloat
% \renewcommand{\subfloat}[2][\relax]{\MYorigsubfloat[]{#2}}
% However, the above trick will not work if both optional arguments of
% the \subfloat command are used. Furthermore, there needs to be a
% description of each subfigure *somewhere* and endfloat does not add
% subfigure captions to its list of figures. Thus, the best approach is to
% avoid the use of subfigure captions (many IEEE journals avoid them anyway)
% and instead reference/explain all the subfigures within the main caption.
% The latest version of endfloat.sty and its documentation can obtained at:
% http://www.ctan.org/pkg/endfloat
%
% The IEEEtran \ifCLASSOPTIONcaptionsoff conditional can also be used
% later in the document, say, to conditionally put the References on a 
% page by themselves.




% *** PDF, URL AND HYPERLINK PACKAGES ***
%
%\usepackage{url}
% url.sty was written by Donald Arseneau. It provides better support for
% handling and breaking URLs. url.sty is already installed on most LaTeX
% systems. The latest version and documentation can be obtained at:
% http://www.ctan.org/pkg/url
% Basically, \url{my_url_here}.




% *** Do not adjust lengths that control margins, column widths, etc. ***
% *** Do not use packages that alter fonts (such as pslatex).         ***
% There should be no need to do such things with IEEEtran.cls V1.6 and later.
% (Unless specifically asked to do so by the journal or conference you plan
% to submit to, of course. )


% correct bad hyphenation here
\hyphenation{op-tical net-works semi-conduc-tor}


\begin{document}
%
% paper title
% Titles are generally capitalized except for words such as a, an, and, as,
% at, but, by, for, in, nor, of, on, or, the, to and up, which are usually
% not capitalized unless they are the first or last word of the title.
% Linebreaks \\ can be used within to get better formatting as desired.
% Do not put math or special symbols in the title.
\title{ST-JMCT: A Spatio-Temporal Joint-Masked Cross-Modal Transformer for Excavator Video-to-Action Control}
%
%
% author names and IEEE memberships
% note positions of commas and nonbreaking spaces ( ~ ) LaTeX will not break
% a structure at a ~ so this keeps an author's name from being broken across
% two lines.
% use \thanks{} to gain access to the first footnote area
% a separate \thanks must be used for each paragraph as LaTeX2e's \thanks
% was not built to handle multiple paragraphs
%

\author{Michael~Shell,~\IEEEmembership{Member,~IEEE,}
        John~Doe,~\IEEEmembership{Fellow,~OSA,}
        and~Jane~Doe,~\IEEEmembership{Life~Fellow,~IEEE}% <-this % stops a space
\thanks{M. Shell was with the Department
of Electrical and Computer Engineering, Georgia Institute of Technology, Atlanta,
GA, 30332 USA e-mail: (see http://www.michaelshell.org/contact.html).}% <-this % stops a space
\thanks{J. Doe and J. Doe are with Anonymous University.}% <-this % stops a space
\thanks{Manuscript received April 19, 2005; revised August 26, 2015.}}

% note the % following the last \IEEEmembership and also \thanks - 
% these prevent an unwanted space from occurring between the last author name
% and the end of the author line. i.e., if you had this:
% 
% \author{....lastname \thanks{...} \thanks{...} }
%                     ^------------^------------^----Do not want these spaces!
%
% a space would be appended to the last name and could cause every name on that
% line to be shifted left slightly. This is one of those "LaTeX things". For
% instance, "\textbf{A} \textbf{B}" will typeset as "A B" not "AB". To get
% "AB" then you have to do: "\textbf{A}\textbf{B}"
% \thanks is no different in this regard, so shield the last } of each \thanks
% that ends a line with a % and do not let a space in before the next \thanks.
% Spaces after \IEEEmembership other than the last one are OK (and needed) as
% you are supposed to have spaces between the names. For what it is worth,
% this is a minor point as most people would not even notice if the said evil
% space somehow managed to creep in.



% The paper headers
\markboth{Journal of \LaTeX\ Class Files,~Vol.~14, No.~8, August~2015}%
{Shell \MakeLowercase{\textit{et al.}}: Bare Demo of IEEEtran.cls for IEEE Journals}
% The only time the second header will appear is for the odd numbered pages
% after the title page when using the twoside option.
% 
% *** Note that you probably will NOT want to include the author's ***
% *** name in the headers of peer review papers.                   ***
% You can use \ifCLASSOPTIONpeerreview for conditional compilation here if
% you desire.




% If you want to put a publisher's ID mark on the page you can do it like
% this:
%\IEEEpubid{0000--0000/00\$00.00~\copyright~2015 IEEE}
% Remember, if you use this you must call \IEEEpubidadjcol in the second
% column for its text to clear the IEEEpubid mark.



% use for special paper notices
%\IEEEspecialpapernotice{(Invited Paper)}




% make the title area
\maketitle

% As a general rule, do not put math, special symbols or citations
% in the abstract or keywords.
\begin{abstract}
Autonomous excavation requires models that infer coordinated joint motion from
dynamic visual observations without relying on measured machine states. We
present ST-JMCT, a spatio-temporal joint-masked cross-modal transformer that
predicts next-step boom, arm, bucket, and swing configurations from RGB video
and LiDAR-derived elevation maps. Independent visual encoders and bidirectional
cross-modal attention combine appearance and geometric information. Eight
joint-specific spatio-temporal masks, one for each modality--joint pair, then
guide the action decoder, emphasizing local evidence for the working equipment
while retaining global context for swing estimation. A within-modality
diversity objective prevents different joint masks from collapsing onto the
same region while permitting corresponding RGB and elevation masks to agree.
Periodic angle representations and training-only pose and swing-velocity heads
provide auxiliary supervision without introducing joint-state inputs at
inference. We evaluate ST-JMCT on the public EXC Motion dataset released by
Fuxi Robot, which contains operating sequences from three excavator types. The
proposed framework provides an interpretable multimodal video-to-action
formulation for excavator automation.
\end{abstract}

% Note that keywords are not normally used for peerreview papers.
\begin{IEEEkeywords}
autonomous excavators, cross-modal attention, joint-angle prediction,
multimodal learning, spatio-temporal attention, video-to-action.
\end{IEEEkeywords}






% For peer review papers, you can put extra information on the cover
% page as needed:
% \ifCLASSOPTIONpeerreview
% \begin{center} \bfseries EDICS Category: 3-BBND \end{center}
% \fi
%
% For peerreview papers, this IEEEtran command inserts a page break and
% creates the second title. It will be ignored for other modes.
\IEEEpeerreviewmaketitle



\section{Introduction}
\label{sec:introduction}
% The Introduction will be completed after the Methodology section.
\iffalse
% The very first letter is a 2 line initial drop letter followed
% by the rest of the first word in caps.
% 
% form to use if the first word consists of a single letter:
% \IEEEPARstart{A}{demo} file is ....
% 
% form to use if you need the single drop letter followed by
% normal text (unknown if ever used by the IEEE):
% \IEEEPARstart{A}{}demo file is ....
% 
% Some journals put the first two words in caps:
% \IEEEPARstart{T}{his demo} file is ....
% 
% Here we have the typical use of a "T" for an initial drop letter
% and "HIS" in caps to complete the first word.
\IEEEPARstart{T}{his} demo file is intended to serve as a ``starter file''
for IEEE journal papers produced under \LaTeX\ using
IEEEtran.cls version 1.8b and later.
% You must have at least 2 lines in the paragraph with the drop letter
% (should never be an issue)
I wish you the best of success.

\hfill mds
 
\hfill August 26, 2015

\subsection{Subsection Heading Here}
Subsection text here.

% needed in second column of first page if using \IEEEpubid
%\IEEEpubidadjcol

\subsubsection{Subsubsection Heading Here}
Subsubsection text here.


% An example of a floating figure using the graphicx package.
% Note that \label must occur AFTER (or within) \caption.
% For figures, \caption should occur after the \includegraphics.
% Note that IEEEtran v1.7 and later has special internal code that
% is designed to preserve the operation of \label within \caption
% even when the captionsoff option is in effect. However, because
% of issues like this, it may be the safest practice to put all your
% \label just after \caption rather than within \caption{}.
%
% Reminder: the "draftcls" or "draftclsnofoot", not "draft", class
% option should be used if it is desired that the figures are to be
% displayed while in draft mode.
%
%\begin{figure}[!t]
%\centering
%\includegraphics[width=2.5in]{myfigure}
% where an .eps filename suffix will be assumed under latex, 
% and a .pdf suffix will be assumed for pdflatex; or what has been declared
% via \DeclareGraphicsExtensions.
%\caption{Simulation results for the network.}
%\label{fig_sim}
%\end{figure}

% Note that the IEEE typically puts floats only at the top, even when this
% results in a large percentage of a column being occupied by floats.


% An example of a double column floating figure using two subfigures.
% (The subfig.sty package must be loaded for this to work.)
% The subfigure \label commands are set within each subfloat command,
% and the \label for the overall figure must come after \caption.
% \hfil is used as a separator to get equal spacing.
% Watch out that the combined width of all the subfigures on a 
% line do not exceed the text width or a line break will occur.
%
%\begin{figure*}[!t]
%\centering
%\subfloat[Case I]{\includegraphics[width=2.5in]{box}%
%\label{fig_first_case}}
%\hfil
%\subfloat[Case II]{\includegraphics[width=2.5in]{box}%
%\label{fig_second_case}}
%\caption{Simulation results for the network.}
%\label{fig_sim}
%\end{figure*}
%
% Note that often IEEE papers with subfigures do not employ subfigure
% captions (using the optional argument to \subfloat[]), but instead will
% reference/describe all of them (a), (b), etc., within the main caption.
% Be aware that for subfig.sty to generate the (a), (b), etc., subfigure
% labels, the optional argument to \subfloat must be present. If a
% subcaption is not desired, just leave its contents blank,
% e.g., \subfloat[].


% An example of a floating table. Note that, for IEEE style tables, the
% \caption command should come BEFORE the table and, given that table
% captions serve much like titles, are usually capitalized except for words
% such as a, an, and, as, at, but, by, for, in, nor, of, on, or, the, to
% and up, which are usually not capitalized unless they are the first or
% last word of the caption. Table text will default to \footnotesize as
% the IEEE normally uses this smaller font for tables.
% The \label must come after \caption as always.
%
%\begin{table}[!t]
%% increase table row spacing, adjust to taste
%\renewcommand{\arraystretch}{1.3}
% if using array.sty, it might be a good idea to tweak the value of
% \extrarowheight as needed to properly center the text within the cells
%\caption{An Example of a Table}
%\label{table_example}
%\centering
%% Some packages, such as MDW tools, offer better commands for making tables
%% than the plain LaTeX2e tabular which is used here.
%\begin{tabular}{|c||c|}
%\hline
%One & Two\\
%\hline
%Three & Four\\
%\hline
%\end{tabular}
%\end{table}


% Note that the IEEE does not put floats in the very first column
% - or typically anywhere on the first page for that matter. Also,
% in-text middle ("here") positioning is typically not used, but it
% is allowed and encouraged for Computer Society conferences (but
% not Computer Society journals). Most IEEE journals/conferences use
% top floats exclusively. 
% Note that, LaTeX2e, unlike IEEE journals/conferences, places
% footnotes above bottom floats. This can be corrected via the
% \fnbelowfloat command of the stfloats package.

\fi

\section{Methodology}
\label{sec:method}

\subsection{Problem Formulation and Architecture Overview}
\label{subsec:overview}

We formulate excavator video-to-action learning as next-step joint-state
prediction from a fixed window of visual observations. Let
$\mathbf{X}^{r}_{t-T+1:t}$ and $\mathbf{X}^{e}_{t-T+1:t}$ denote an RGB clip
and its synchronized LiDAR-derived elevation clip, respectively, where
$\mathbf{X}^{m}\in\mathbb{R}^{B\times T\times3\times H\times W}$, $m\in
\{r,e\}$, and $T=8$. A categorical machine identifier $c$ distinguishes the
75, 306, and 490 excavator types. ST-JMCT predicts the next configuration of
the boom, arm, bucket, and swing joints:
\begin{equation}
\widehat{\mathbf{y}}_{t+1}=
F_{\Theta}\!\left(\mathbf{X}^{r}_{t-T+1:t},
\mathbf{X}^{e}_{t-T+1:t},c\right),
\label{eq:task}
\end{equation}
where $\widehat{\mathbf{y}}_{t+1}\in\mathbb{R}^{B\times8}$ contains one
sine--cosine pair for each joint. Measured joint positions provide targets and
auxiliary supervision during training but are not consumed by the normal
inference path.

Fig.~\ref{fig:framework} summarizes the architecture. The two visual streams
first extract multi-scale appearance and geometric features. Bidirectional
cross-modal attention exchanges information between the streams, after which
a shared TemporalMaskMixer is applied separately to each modality. Four
joint-specific mask heads produce one mask per joint and modality. The
temporally enhanced tokens are then fused, encoded as a global memory, and
read by four mask-guided joint queries. Finally, machine-specific action heads
map the decoded joint features to the next joint configuration.

\begin{figure*}[!t]
    \centering
    \includegraphics[width=\textwidth]{Figure1.jpg}
    \caption{Overview of ST-JMCT. Synchronized RGB and LiDAR-derived elevation
    clips are encoded independently, exchanged through bidirectional
    cross-modal attention, and mixed temporally before joint-specific masks
    are generated. The masks guide four joint queries, while training-only
    auxiliary heads provide pose and swing-velocity supervision.}
    \label{fig:framework}
\end{figure*}

\subsection{Dual-Stream Spatio-Temporal Representation}
\label{subsec:dual_stream}

RGB observations describe machine appearance, scene layout, and visible
motion, whereas elevation maps provide geometric evidence about the machine
and surrounding terrain. Directly sharing the complete feature extractor
would force both modalities into the same low-level statistics. We therefore
use two CSPDarknet-style backbones followed by independent FPN--PAN necks. A
lightweight elevation adapter, composed of two $3\times3$ convolutions with
batch normalization and a SiLU activation, maps the rendered elevation input
to the domain expected by its visual backbone.

For each frame, the two branches produce three pyramid features at spatial
resolutions $28\times28$, $14\times14$, and $7\times7$. A spatial grid head
resamples the three scales to a common $G\times G$ grid, with $G=14$,
concatenates them, and projects each modality to $D=512$ dimensions. This
gives $\mathbf{Z}^{r},\mathbf{Z}^{e}\in
\mathbb{R}^{B\times T\times G\times G\times D}$. Before temporal mixing, the
two streams exchange complementary evidence using bidirectional multi-head
cross-attention. After flattening the temporal-spatial axes, the interaction is
\begin{align}
\overline{\mathbf{Z}}^{r} &={\rm MHA}
\left(\mathbf{Z}^{r},\mathbf{Z}^{e},\mathbf{Z}^{e}\right), \\
\overline{\mathbf{Z}}^{e} &={\rm MHA}
\left(\mathbf{Z}^{e},\mathbf{Z}^{r},\mathbf{Z}^{r}\right),
\label{eq:cross_modal}
\end{align}
where the three arguments denote query, key, and value tokens.

The cross-modal outputs are subsequently processed by the TemporalMaskMixer.
For every grid location $p$, tokens from the same spatial position across the
$T$ frames are reshaped to a sequence in
$\mathbb{R}^{T\times D}$ and processed by a Transformer encoder:
\begin{equation}
\widetilde{\mathbf{Z}}^{m}_{:,p}=
\mathcal{T}\!\left(\overline{\mathbf{Z}}^{m}_{:,p}\right),
\qquad m\in\{r,e\}.
\label{eq:temporal_mixer}
\end{equation}
The same temporal module is applied separately to the RGB and elevation
streams. Consequently, the mask generators observe both cross-modal evidence
and motion over the input window while the modality-specific token sets remain
available.

\subsection{Joint-Specific Cross-Modal Mask Generation}
\label{subsec:joint_masks}

Different excavator joints depend on different visual regions and spatial
scales. The bucket is closely associated with the end effector and the terrain
contact region, whereas swing estimation requires a wider reference covering
the machine orientation and surrounding scene. A single shared attention map
cannot express these distinct requirements. ST-JMCT therefore uses four
independent mask heads conditioned on learnable joint embeddings. The same
joint head is applied to each modality so that its semantic role is consistent
between RGB and elevation features.

Let $\mathbf{e}_{j}\in\mathbb{R}^{D}$ be the embedding of joint
$j\in\mathcal{J}=\{\text{boom},\text{arm},\text{bucket},\text{swing}\}$.
For modality $m$, the corresponding mask is
\begin{equation}
\mathbf{M}^{m}_{j}=
\sigma\!\left(
\mathbf{W}_{2,j}\,\mathrm{GELU}\!\left[
\mathbf{W}_{1,j}\left(\widetilde{\mathbf{Z}}^{m}+\mathbf{e}_{j}\right)
\right]\right),
\label{eq:mask_generation}
\end{equation}
where $\sigma$ is the sigmoid function. The complete mask tensor has shape
$\mathbf{M}\in\mathbb{R}^{B\times2\times4\times T\times G\times G}$,
corresponding to two modalities and four joints. For joint $j$, the RGB and
elevation masks are combined by a differentiable soft union,
\begin{equation}
\widetilde{\mathbf{M}}_{j}=
1-\left(1-\mathbf{M}^{r}_{j}\right)
  \left(1-\mathbf{M}^{e}_{j}\right),
\label{eq:mask_union}
\end{equation}
which preserves a high response when either modality identifies a relevant
region.

To prevent different joint masks from converging to an identical solution,
we regularize different-joint pairs within each modality. After flattening
and $\ell_2$-normalizing every temporal-spatial mask, the diversity term is
\begin{equation}
\mathcal{L}_{\mathrm{div}}=
\frac{1}{2}\underset{b,m,j\neq k}{\operatorname{mean}}
\left[\max\!\left(0,
\left\langle\bar{\mathbf{M}}^{m}_{b,j},
\bar{\mathbf{M}}^{m}_{b,k}\right\rangle-\delta\right)^{2}\right],
\label{eq:diversity}
\end{equation}
with margin $\delta=0.5$. Same-joint RGB/elevation pairs are excluded from
this penalty because cross-modal agreement may correspond to the same physical
machine component or terrain region.

\subsection{Mask-Guided Joint Decoding}
\label{subsec:decoder}

The temporally enhanced RGB and elevation features are concatenated and
linearly projected back to $D$ dimensions. Fixed two-dimensional spatial
encodings, learned temporal encodings, and a learned embedding for machine
type $c$ are added to the fused tokens. A three-layer bidirectional Transformer
encoder then produces the visual memory
$\mathbf{H}\in\mathbb{R}^{B\times(TG^{2})\times D}$. The encoder itself is
not hard-gated by the masks, because suppressing tokens at this stage would
remove global context needed by the swing joint.

Four learned queries read the shared memory through mask-guided cross-attention.
For query $\boldsymbol{\phi}_{j}$, the joint mask affects both the attention
logits and the value features:
\begin{align}
\mathbf{A}_{j} &= \operatorname{softmax}\!\left(
\frac{\mathbf{Q}_{j}\mathbf{K}^{\mathsf{T}}}{\sqrt{d_h}}+
\lambda_{m,j}\log\!\left(\widetilde{\mathbf{M}}_{j}+\epsilon\right)
\right), \\
\mathbf{V}'_{j} &= \mathbf{V}\odot
\left(1+\lambda_{v,j}\widetilde{\mathbf{M}}_{j}\right), \\
\mathbf{h}_{j} &= \mathbf{A}_{j}\mathbf{V}'_{j},
\label{eq:mask_attention}
\end{align}
where $d_h$ is the attention-head dimension. The residual value modulation
retains the unmasked visual signal and amplifies mask-responsive locations,
instead of zeroing all features outside the selected region.

The mask strengths reflect the visual scope of each joint, as summarized in
Table~\ref{tab:mask_strength}. Boom, arm, and bucket use progressively stronger
logit guidance. Swing retains global attention by setting
$\lambda_{m,\mathrm{swing}}=0$, while its mask still provides weak residual
value enhancement.

\begin{table}[!t]
\caption{Joint-Specific Mask Guidance}
\label{tab:mask_strength}
\centering
\begin{tabular}{lccc}
\toprule
Joint & $\lambda_m$ & $\lambda_v$ & Visual scope \\
\midrule
Boom   & 1.5 & 0.8 & Semi-local \\
Arm    & 2.0 & 1.0 & Local \\
Bucket & 2.5 & 1.0 & Strong local \\
Swing  & 0.0 & 0.5 & Global \\
\bottomrule
\end{tabular}
\end{table}

\subsection{Action Prediction and Training Objectives}
\label{subsec:objectives}

The decoded features are passed to per-machine, per-joint prediction heads.
For a sample with machine identifier $c$, only the four heads associated with
that machine are selected. Each head predicts a pair
$\widehat{\mathbf{y}}_{j}=[\hat{s}_{j},\hat{c}_{j}]$, and the corresponding
angle is recovered as
\begin{equation}
\hat{q}_{t+1,j}=\operatorname{atan2}(\hat{s}_{j},\hat{c}_{j}).
\label{eq:angle_decode}
\end{equation}
This periodic representation avoids a discontinuous regression target at the
$-\pi/\pi$ boundary.

The main prediction objective is a weighted mean-squared error over the four
sine--cosine pairs,
\begin{equation}
\mathcal{L}_{\mathrm{act}}=
\frac{\sum_{j\in\mathcal{J}}w_j
\left\|\widehat{\mathbf{y}}_{j}-\mathbf{y}_{j}\right\|_2^2}
{\sum_{j\in\mathcal{J}}w_j},
\label{eq:action_loss}
\end{equation}
where $w_{\mathrm{swing}}=2$ and the other joint weights are one. A unit-circle
penalty constrains each predicted pair,
$\mathcal{L}_{\mathrm{circ}}=\operatorname{mean}_{j}
(\hat{s}_{j}^{2}+\hat{c}_{j}^{2}-1)^{2}$.

Mask learning uses three additional terms. An area loss keeps the boom, arm,
and bucket masks within an average active-area interval of $[0.05,0.30]$ and
the broader swing mask within $[0.10,0.70]$. The diversity loss in
Eq.~\eqref{eq:diversity} separates different joints within each modality, and
a temporal smoothness loss penalizes abrupt changes between adjacent masks.
The regularization weight is gradually reduced after epoch 40 so that the
later optimization is dominated by action prediction.

Two training-only heads provide additional motion supervision. The pose head
predicts the sine--cosine representation of the current four-joint pose from
the last-frame visual memory. A separate swing-velocity head predicts the
wrapped angular increment between the current and next swing angles. The full
training objective is
\begin{equation}
\begin{split}
\mathcal{L}={}&\mathcal{L}_{\mathrm{act}}
+0.1\mathcal{L}_{\mathrm{circ}}
+\rho(e)\left(\mathcal{L}_{\mathrm{area}}
+\mathcal{L}_{\mathrm{div}}+\mathcal{L}_{\mathrm{temp}}\right)\\
&+\lambda_{p}\mathcal{L}_{\mathrm{pose}}
+\lambda_{v}\mathcal{L}_{\mathrm{vel}},
\end{split}
\label{eq:total_loss}
\end{equation}
where $\lambda_p=0.05$, $\lambda_v=0.10$, and $\rho(e)$ denotes the
epoch-dependent mask-regularization scale. Both auxiliary heads are discarded
during inference; thus, predicted actions depend only on the two visual clips
and the machine identifier.


\section{Experiments and Evaluation}
\label{sec:experiments}

\subsection{Dataset and Task Configuration}
\label{subsec:dataset}

We evaluated ST-JMCT on the public Excavator Motion dataset released by Fuxi
Robot. The dataset contains 30 excavation episodes collected from three
hydraulic excavators, denoted by their model identifiers 75, 306, and 490.
Each model contributes ten episodes, and only these three excavator types are
used in the experiments. Every episode provides synchronized RGB images,
LiDAR-derived elevation images, and four-dimensional joint positions ordered
as boom, arm, bucket, and swing.

The episodes were split independently within each excavator type using a fixed
random seed of 42. Eight episodes from each machine were assigned to training
and the remaining two to validation, resulting in 24 training episodes and six
held-out validation episodes (Table~\ref{tab:dataset_split}). Splitting at the
episode level prevents temporally adjacent windows from the same trajectory
from appearing in both subsets. The current protocol does not use an
independent test subset; a separate test split will be fixed before the final
comparative evaluation.

\begin{table}[!t]
\caption{Episode-Level Dataset Split}
\label{tab:dataset_split}
\centering
\begin{tabular}{lccc}
\toprule
Excavator & Total episodes & Training & Validation \\
\midrule
75  & 10 & 8 & 2 \\
306 & 10 & 8 & 2 \\
490 & 10 & 8 & 2 \\
\midrule
Total & 30 & 24 & 6 \\
\bottomrule
\end{tabular}
\end{table}

Sliding windows were extracted with a stride of one frame. Each sample uses
$T=8$ synchronized RGB and elevation frames to predict the four joint angles
at the next frame. Both image modalities were resized to $224\times224$,
converted to RGB ordering, scaled to $[0,1]$, and normalized using ImageNet
statistics. Ground-truth angles were converted to four sine--cosine pairs for
optimization, while reported errors were computed after recovering angles with
$\operatorname{atan2}$.

\subsection{Experimental Setup}
\label{subsec:experimental_setup}

Unless otherwise stated, all models used a hidden dimension of 512, eight
input frames, a $14\times14$ token grid, and three layers in the global
Transformer encoder. ST-JMCT was trained for 80 epochs with a batch size of 8
using AdamW. The initial learning rate was $3\times10^{-5}$ and the weight
decay was $10^{-3}$. A linear warm-up covered the first 5\% of optimization
steps and was followed by cosine annealing. Gradients were clipped to a norm of
1.0, and an exponential moving average of the model parameters was maintained
with a decay factor of 0.999. The checkpoint with the lowest validation loss
was used for evaluation. Training was performed on [GPU and software
environment to be specified].

Performance was measured in angular space using mean absolute error (MAE) and
the coefficient of determination ($R^2$). Let $q_{i,j}$ and
$\hat{q}_{i,j}$ denote the target and predicted angles for sample $i$ and joint
$j$. For swing, the residual was wrapped to $[-\pi,\pi]$ as
\begin{equation}
d_{i,j}=\operatorname{atan2}\!\left(
\sin(\hat{q}_{i,j}-q_{i,j}),
\cos(\hat{q}_{i,j}-q_{i,j})\right),
\label{eq:circular_error}
\end{equation}
whereas the direct angular difference was used for the remaining joints. The
per-joint MAE is
\begin{equation}
\operatorname{MAE}_{j}=\frac{1}{N}\sum_{i=1}^{N}|d_{i,j}|.
\label{eq:mae}
\end{equation}
For swing, $R^2$ was computed using wrapped deviations from the circular mean;
the conventional linear definition was used for boom, arm, and bucket. We
report overall averages together with joint-wise and excavator-wise values so
that performance on a single joint or machine cannot dominate the aggregate
score.

\subsection{Comparison with a Standard Transformer}
\label{subsec:baseline}

Reliable published baselines for excavator video-to-action prediction under
the same sensing and evaluation protocol are not currently available. We
therefore use a standard Transformer as a controlled reference rather than
claiming comparison with the state of the art. The baseline receives the same
RGB and elevation clips and machine identifiers, uses the same input
preprocessing and sine--cosine targets, and is trained with the same data split
and optimization schedule. It encodes the dual-stream visual features with a
conventional temporal Transformer and predicts the next joint configuration
without bidirectional cross-modal token interaction, TemporalMaskMixer, or
joint-specific mask-guided decoding.

Table~\ref{tab:main_results} will report the final comparison. In addition to
prediction accuracy, we will include the number of trainable parameters and
inference latency to distinguish architectural gains from increased
computational cost. Numerical statements will be added only after both models
have completed evaluation under the fixed protocol.

\begin{table}[!t]
\caption{Comparison with the Standard Transformer under the Same Data Split
and Training Protocol. Lower MAE and latency are better; higher $R^2$ is
better.}
\label{tab:main_results}
\centering
\resizebox{\columnwidth}{!}{%
\begin{tabular}{lccccc}
\toprule
Model & Parameters (M) & MAE (rad) $\downarrow$ & $R^2$ $\uparrow$
& Latency (ms) $\downarrow$ & Input modalities \\
\midrule
Standard Transformer & -- & -- & -- & -- & RGB + elevation \\
ST-JMCT              & -- & -- & -- & -- & RGB + elevation \\
\bottomrule
\end{tabular}
}
\end{table}

\subsection{Ablation Studies}
\label{subsec:ablation}

The ablation study is designed to associate each architectural contribution
with a measurable change in prediction error. Starting from the complete
ST-JMCT model, we remove the elevation stream, bidirectional cross-modal
interaction, TemporalMaskMixer, and joint-specific mask guidance in turn. A
shared-mask variant replaces the four joint heads with one common mask, testing
whether joint-conditioned spatial selection is more useful than generic visual
attention. Further variants remove the within-modality mask-diversity loss or
the two auxiliary objectives. All variants retain the same episode split,
training schedule, output representation, and model-selection criterion.

The primary ablations and their joint-wise errors will be summarized in
Table~\ref{tab:ablation_results}. The cross-modal and temporal modules validate
multimodal and motion modeling, respectively, whereas the no-mask and
shared-mask variants isolate the contribution of joint-specific spatial
guidance. The auxiliary-loss ablation tests whether the training-only signals
improve action prediction without increasing inference-time input
requirements.

\begin{table}[!t]
\caption{Ablation Study of ST-JMCT. Values are angular MAE in radians; lower is
better.}
\label{tab:ablation_results}
\centering
\resizebox{\columnwidth}{!}{%
\begin{tabular}{lccccc}
\toprule
Variant & Boom & Arm & Bucket & Swing & Mean \\
\midrule
Full ST-JMCT                    & -- & -- & -- & -- & -- \\
Without elevation stream       & -- & -- & -- & -- & -- \\
Without cross-modal interaction& -- & -- & -- & -- & -- \\
Without TemporalMaskMixer      & -- & -- & -- & -- & -- \\
Without joint-specific masks   & -- & -- & -- & -- & -- \\
Shared mask for all joints     & -- & -- & -- & -- & -- \\
Without mask-diversity loss    & -- & -- & -- & -- & -- \\
Without auxiliary objectives   & -- & -- & -- & -- & -- \\
\bottomrule
\end{tabular}
}
\end{table}

\subsection{Joint-, Machine-Wise, and Qualitative Analysis}
\label{subsec:detailed_analysis}

Aggregate metrics alone may conceal differences between joint dynamics and
excavator platforms. We therefore report MAE and $R^2$ separately for boom,
arm, bucket, and swing, and repeat the analysis for the 75, 306, and 490
machines. Particular attention is given to swing because its periodic motion
requires wrap-aware evaluation and receives a larger weight in the training
objective. Episode-wise results will additionally be inspected to identify
whether isolated trajectories account for a disproportionate fraction of the
error.

The quantitative breakdown will be complemented by two forms of qualitative
evidence. First, predicted and ground-truth joint-angle trajectories will be
overlaid for representative episodes from all three excavators. Second, the
RGB and elevation masks for each joint will be visualized on the corresponding
frames. These visualizations will be used to examine whether boom, arm, and
bucket queries concentrate on local machine structures and interaction regions
while the swing query preserves broader scene context. We will also include
representative failure cases, particularly rapid transitions, partial
occlusion, and visually ambiguous motion, to delimit the conditions under
which the model remains reliable.




\section{Conclusion}
The conclusion goes here.





% if have a single appendix:
%\appendix[Proof of the Zonklar Equations]
% or
%\appendix  % for no appendix heading
% do not use \section anymore after \appendix, only \section*
% is possibly needed

% use appendices with more than one appendix
% then use \section to start each appendix
% you must declare a \section before using any
% \subsection or using \label (\appendices by itself
% starts a section numbered zero.)
%


\appendices
\section{Proof of the First Zonklar Equation}
Appendix one text goes here.

% you can choose not to have a title for an appendix
% if you want by leaving the argument blank
\section{}
Appendix two text goes here.


% use section* for acknowledgment
\section*{Acknowledgment}


The authors would like to thank...


% Can use something like this to put references on a page
% by themselves when using endfloat and the captionsoff option.
\ifCLASSOPTIONcaptionsoff
  \newpage
\fi



% trigger a \newpage just before the given reference
% number - used to balance the columns on the last page
% adjust value as needed - may need to be readjusted if
% the document is modified later
%\IEEEtriggeratref{8}
% The "triggered" command can be changed if desired:
%\IEEEtriggercmd{\enlargethispage{-5in}}

% references section

% can use a bibliography generated by BibTeX as a .bbl file
% BibTeX documentation can be easily obtained at:
% http://mirror.ctan.org/biblio/bibtex/contrib/doc/
% The IEEEtran BibTeX style support page is at:
% http://www.michaelshell.org/tex/ieeetran/bibtex/
%\bibliographystyle{IEEEtran}
% argument is your BibTeX string definitions and bibliography database(s)
%\bibliography{IEEEabrv,../bib/paper}
%
% <OR> manually copy in the resultant .bbl file
% set second argument of \begin to the number of references
% (used to reserve space for the reference number labels box)
\begin{thebibliography}{1}

\bibitem{IEEEhowto:kopka}
H.~Kopka and P.~W. Daly, \emph{A Guide to \LaTeX}, 3rd~ed.\hskip 1em plus
  0.5em minus 0.4em\relax Harlow, England: Addison-Wesley, 1999.

\end{thebibliography}

% biography section
% 
% If you have an EPS/PDF photo (graphicx package needed) extra braces are
% needed around the contents of the optional argument to biography to prevent
% the LaTeX parser from getting confused when it sees the complicated
% \includegraphics command within an optional argument. (You could create
% your own custom macro containing the \includegraphics command to make things
% simpler here.)
%\begin{IEEEbiography}[{\includegraphics[width=1in,height=1.25in,clip,keepaspectratio]{mshell}}]{Michael Shell}
% or if you just want to reserve a space for a photo:

\begin{IEEEbiography}{Michael Shell}
Biography text here.
\end{IEEEbiography}

% if you will not have a photo at all:
\begin{IEEEbiographynophoto}{John Doe}
Biography text here.
\end{IEEEbiographynophoto}

% insert where needed to balance the two columns on the last page with
% biographies
%\newpage

\begin{IEEEbiographynophoto}{Jane Doe}
Biography text here.
\end{IEEEbiographynophoto}

% You can push biographies down or up by placing
% a \vfill before or after them. The appropriate
% use of \vfill depends on what kind of text is
% on the last page and whether or not the columns
% are being equalized.

%\vfill

% Can be used to pull up biographies so that the bottom of the last one
% is flush with the other column.
%\enlargethispage{-5in}



% that's all folks
\end{document}
